\SourceLeaf{243}%
bis ich mich einmal mit Musse wieder in dieses Fach werfen kann, und das Lesen russischer Bücher mir noch geläufiger ist als jetzt. Irre ich nicht, so ist auch ein Aufsatz des p. \textsc{Lobatschewsky}, vielleicht eine Übersetzung aus den \gausscyr{Записки}, in \textsc{Crelles} Journal, was ich aber in diesem Augenblick nicht nachsehen kann. \mbox{.\ .\ .\ .\ .\ .\ .}

\par\bigskip\Needspace{6\baselineskip}

\begin{center}\textsc{Gauss} an \textsc{Gerling}. Göttingen, 8. Februar 1844.\end{center}

Es wird Ihnen, mein theuerster Freund, vielleicht nicht unlieb sein, wenn ich zu den literarischen Notizen, welche ich in meinem letzten Briefe mittheilte, noch eine oder die andere hinzufüge: in die Sache selbst kann ich freilich jetzt nicht tiefer eingehen.

\textsc{Lobatschewskys} Aufsatz in \textsc{Crelles} Journal steht Band 17 pag. 295 ff. Ich finde, dass derselbe nur eine freie Übersetzung des russischen Aufsatzes [\gausscyr{Воображаемая геометрія}] im Jahrgang 1835 der Gelehrten Schriften der Kasanschen Universität ist, wo man eben da auch anstossen wird, wo diess in dem deutschen Aufsatze der Fall ist. In diesem stossen Sie an S. 296 Zeile 10 bei den Worten \textls{J’ai démontré} \emph{etc.}, womit dem Leser, der weiter nichts hat wie diesen Aufsatz, wenig gedient ist. Ebenso S. 303 oben \textls{J’ai prouvé ailleurs} \emph{etc.}, wozu man dieselbe Bemerkung machen muss. Der frühere Aufsatz, worauf sich diess zu beziehen \textls{scheint}, wird wohl derselbe sein, der in einer Note des erwähnten russischen Aufsatzes angeführt wird als unter dem Titel: Über die Anfänge oder Principe der Geometrie stehend in \gausscyr{Казанскомъ Вѣстникѣ} (Kasanschen Boten) für 1829 und 1830. Zugleich wird dabei bemerkt, dass eine sehr kränkende Kritik dieser Abhandlung in No. 41 eines andern russischen, wie ich vermuthe in Petersburg erscheinenden, Journals »Sohn des Vaterlandes« \gausscyr{Сынъ Отечества} von 1834 stehe, wogegen \textsc{Lobatschewsky} eine Antikritik eingeschickt habe, die aber bis Anfang 1835 nicht aufgenommen sei.

Mit diesen literarischen Notizen ist uns nun freilich auch wenig geholfen, da in Deutschland schwerlich ein Exemplar des Kasanschen Boten von 1829---%
